
\(U\) \uline{de} \(x\), \uline{et une suite de sections} \(f_1,\ldots,f_p,g_1,\ldots,g_r,h_1,\ldots,h_s\) \uline{de} \(\underline{O}_U\) \uline{telles que le complexe de Koszul associé à} \(f_1,\ldots,f_p,g_1,\ldots,g_r\) \uline{soit une résolution de} \(\underline{O}_{Y\cap U}\), \uline{le complexe de Koszul associé à} \(f_1,\ldots,f_p,h_1,\ldots,h_s\) \uline{une résolution de} \(\underline{O}_{Z\cap U}\), \uline{et le complexe de Koszul associé à} \(f_1,\ldots,f_p,g_1,\ldots,g_r,h_1,\ldots,h_s\) \uline{une résolution de} \(\underline{O}_{T\cap U}\) (\uline{conditions satisfaites en particulier dans le cas noethérien si la suite} \(f_1,\ldots,f_p,g_1,\ldots,g_r,h_1,\ldots,h_s\) \uline{est régulière et si} \(Y=V(f_1,\ldots,f_p,g_1,\ldots,g_r)\), \(Z=V(f_1,\ldots,f_p,h_1,\ldots,h_s)\)). \uline{Soient} \(\underline{I},\underline{J}\) \uline{les idéaux de} \(Y,Z\). \uline{On a alors les isomorphismes suivants}:
\[
\begin{array}{ll}
\text{i)}& \forall i\ge0,\quad
\underline{\Tor}_i^{\underline{O}_X}(\underline{O}_Y,\underline{O}_Z)\xrightarrow{\sim}
\bigwedge^i\underline{\Tor}_1^{\underline{O}_X}(\underline{O}_Y,\underline{O}_Z);\\[0.8ex]
\text{ii)}&
\underline{\Tor}_1^{\underline{O}_X}(\underline{O}_Y,\underline{O}_Z)\simeq
\underline{I}\cap \underline{J}/\underline{I}.\underline{J}\simeq
\Ker\bigl(i^*(\underline{N}_{Y/X})\times j^*(\underline{N}_{Z/X})\to \underline{N}_{T/X}\bigr),
\end{array}
\]
\uline{où on note} \(i\) \uline{l'immersion de} \(T\) \uline{dans} \(Y\) \uline{et} \(j\) \uline{l'immersion de} \(T\) \uline{dans} \(Z\).

\uline{En particulier}, \uline{pour toute immersion fermée régulière} \(Y\to X\), \uline{on a}
\[
\begin{array}{ll}
\text{i')}& \forall i\ge0,\quad
\underline{\Tor}_i^{\underline{O}_X}(\underline{O}_Y,\underline{O}_Y)\xrightarrow{\sim}
\bigwedge^i\underline{\Tor}_1^{\underline{O}_X}(\underline{O}_Y,\underline{O}_Y);\\[0.8ex]
\text{ii')}&
\underline{\Tor}_1^{\underline{O}_X}(\underline{O}_Y,\underline{O}_Y)\simeq \underline{N}_{Y/X}.
\end{array}
\]

i) On a vu en 2.3 que \(\bigoplus_{i\ge0}\underline{\Tor}_i^{\underline{O}_X}(\underline{O}_Y,\underline{O}_Z)\) pouvait être muni d'une structure de \(\underline{O}_T\)-Algèbre; d'après 2.2, cette algèbre est alternée, car l'immersion de \(Y\) (ou \(Z\)) dans \(X\) est régulière. Il existe donc un homomorphisme d'algèbres et un seul:
\begin{equation}\tag{2.5.1}
\bigwedge \underline{\Tor}_1^{\underline{O}_X}(\underline{O}_Y,\underline{O}_Z)\longrightarrow
\bigoplus_{i\ge0}\underline{\Tor}_i^{\underline{O}_X}(\underline{O}_Y,\underline{O}_Z),
\end{equation}
prolongeant l'identité en degré \(1\).

Pour voir que c'est un isomorphisme, on peut le voir localement. On se place donc au voisinage d'un point de \(T\), et on considère des sections \(f_1,\ldots,f_p,g_1,\ldots,g_r,h_1,\ldots,h_s\) de \(\underline{O}_X\) dans ce voisinage, ayant les propriétés indiquées dans l'énoncé de 2.5. Posons alors \(\underline{E}=(\underline{O}_X)^p\), \(\underline{E}'=(\underline{O}_X)^r\), et soient \(u:\underline{E}\to\underline{O}_X\) l'homomorphisme défini par les \(f_i\), \(v:\underline{E}'\to\underline{O}_X\) l'homomorphisme défini par les \(g_j\), \(\underline{K}_\bullet(u)\) et \(\underline{K}_\bullet(v)\) les
